\respcap{cap:uf}
\resp{uf-r1} (a) Manter uma partição de $n$ elementos em conjuntos disjuntos, respondendo se dois elementos estão no mesmo conjunto e unindo conjuntos. (b) \MakeSet$(x)$ cria $\{x\}$ ($pai[x]:=x$); \Find$(x)$ sobe pelos pais até a raiz, que é o representante; \Union$(x,y)$ encontra as duas raízes e faz uma apontar para a outra. (c) O \textbf{pai} de cada elemento na floresta; a raiz guarda também o rank (ou o tamanho). (d) A árvore de menor rank é ligada sob a de maior; o rank só aumenta quando os dois são iguais, e então a árvore pelo menos dobra de tamanho. Por isso uma árvore de rank $h$ tem pelo menos $2^h$ nós, e a altura é no máximo $\log n$ (Teorema~\ref{teo:ufaltura}).

\resp{uf-r2} Ver §\ref{sec:ufexemplos}. (a) Representantes finais $A$ (de $\{A,B,C,D\}$, com rank 2) e $E$ (de $\{E,F\}$): 2 conjuntos. (b) Devolve $A$ e deixa $B,C,D,E$ apontando diretamente para $A$. (c) A raiz é $A$, com rank 2: só no empate a altura cresce. (d) Só com rank, \Find{} custa $\Theta(\log n)$; só com compressão, o total é $O(m\log n)$, com \Find{} individual possivelmente $\Theta(n)$; juntas, $O(m\lgs n)$ (ou $O(m\,\alpha)$).

\resp{uf-l1} A prova é o Teorema~\ref{teo:ufaltura}: por indução, uma árvore de rank (igual à altura, sem compressão) $h$ tem pelo menos $2^h$ nós, pois a altura só aumenta ao unir duas árvores de mesmo rank $h$, cada uma com pelo menos $2^h$ nós. Logo $2^h\le n$. Para atingir a cota com $n=2^k$, una os elementos aos pares (altura 1), depois as árvores resultantes aos pares (altura 2), e assim por diante. Após $k$ rodadas há uma única árvore de altura $k=\log_2n$, que é uma árvore binomial $B_k$.

\resp{uf-l2} Faça \Union$(1,2)$, \Union$(1,3)$, \dots, \Union$(1,n)$, sempre pendurando a raiz da árvore grande sob o novo elemento unitário ($pai[\text{raiz}]:=i$). Forma-se uma cadeia de altura $n-1$ sem nenhum \Find{} caro no caminho; um \Find{} no fundo custa $\Theta(n)$ e achata tudo. Com compressão e união arbitrária, $m$ operações custam $O(m\log n)$. A ideia da prova é classificar cada passo $x\to pai[x]$ da compressão como ``leve'' quando a subárvore de $x$ tem menos da metade dos nós da subárvore do pai: cada nó só passa por $O(\log n)$ passos leves, e cada \Find{} tem $O(\log n)$ passos pesados.

\resp{uf-t4} (a) Union-Find. Como as consultas dominam ($10^8$ contra $10^5$ uniões) e precisam ser rápidas, a \textbf{lista encadeada com ponteiro direto ao representante} responde \Find{} em $O(1)$ \emph{no pior caso}, deixando o custo para as raras uniões: $O(\log n)$ amortizado com ``menor na maior'', no máximo $n\log_2n\approx10^7\cdot23$ atualizações de ponteiro no total. A floresta com rank e compressão também é defensável ($O(\alpha)$ amortizado por operação e menos memória), desde que se diga que a garantia é amortizada; floresta sem balanceamento não serve. (b) \textbf{Anexar o menor ao maior}. Na lista, cada elemento muda de representante no máximo $\log_2n$ vezes, pois o tamanho do seu grupo pelo menos dobra a cada mudança; na árvore, a altura fica no máximo $\log_2n$. Sem isso, listas e árvores degeneram e cada operação pode custar $\Theta(n)$.

\resp{uf-doc} Ver os cenários da §\ref{sec:ufescolha}. (a) Não há como desfazer uniões, e a compressão destrói a forma das árvores. Resolva \emph{offline}, de trás para frente, transformando cada destruição numa união a partir do grafo final; guarde o número de componentes a cada passo e apresente as respostas na ordem original. (b) Floresta sem compressão, só com união por rank ou tamanho (\Find{} em $O(\log n)$), empilhando cada ponteiro e rank alterado por \Union; desfazer é desempilhar e restaurar, em $O(1)$. (c) A floresta só tem ponteiros para cima, e listar membros custaria $O(n)$; use listas encadeadas com ponteiro ao representante e união ``menor na maior'': \emph{MesmoGrupo} em $O(1)$, \emph{ListarMembros} em $O(k)$ e \emph{Juntar} em $O(\log n)$ amortizado.

\resp{uf-cf} Cada bloco de células contíguas é um conjunto, representado por uma \textbf{lista encadeada} dos gatinhos na ordem das células, com cabeça e cauda guardadas no representante. No dia $i$, encontre os representantes de $x_i$ e $y_i$ e concatene uma lista na outra (cauda de uma $\to$ cabeça da outra), unindo também os conjuntos. Como os blocos de $x_i$ e $y_i$ eram vizinhos, colocá-los lado a lado é consistente com alguma ordem inicial. Ao final, percorra a única lista restante. Custo $O(n\,\alpha(n))$.

\resp{uf-sim} \Union$(1,2)$ torna 1 raiz, com rank 1; o mesmo ocorre com 3 (unido a 4), 5 (a 6) e 7 (a 8). \Union$(1,3)$: empate, $pai[3]=1$ e $\rank(1)=2$. \Union$(5,7)$: $pai[7]=5$ e $\rank(5)=2$. \Union$(1,5)$: empate, $pai[5]=1$ e $\rank(1)=3$. Antes do \Find, o vetor $pai$ (posições $1..8$) é $1,1,1,3,1,5,5,7$. \Find$(8)$ percorre $8\to7\to5\to1$, devolve 1 e comprime: $pai=1,1,1,3,1,5,1,1$. Os ranks não mudam: $\rank(1)=3$, $\rank(5)=2$, $\rank(3)=\rank(7)=1$ e os demais são 0.

\resp{uf-ciclo} Faça \MakeSet{} para todo vértice; para cada aresta $uv$, se \Find$(u)=$ \Find$(v)$, a aresta fecha um ciclo; senão, faça \Union$(u,v)$. Custo $O(n+m\,\alpha(n))$. Uma busca em profundidade (Capítulo~\ref{cap:dfs}) também resolve o problema, em $O(n+m)$.

\resp{uf-rankpai} Quando $x$ deixa de ser raiz, ele é ligado a uma raiz $r$ com $\rank(r)>\rank(x)$, ou com $\rank(r)=\rank(x)$, e nesse caso $\rank(r)$ é incrementado; em ambos os casos, $\rank(x)<\rank(r)$. Depois disso, o rank de $x$ nunca muda (só raízes mudam de rank), e o rank de $pai[x]$ nunca diminui. A compressão troca $pai[x]$ por um ancestral, que, pela mesma desigualdade aplicada ao longo do caminho, tem rank ainda maior. Logo a desigualdade se mantém.

\resp{uf-max} Guarde na raiz o maior elemento e o tamanho do conjunto. Na \Union{} de raízes $r$ e $s$, com $s$ indo para baixo de $r$, faça $max[r]:=\max\{max[r],max[s]\}$ e $tam[r]:=tam[r]+tam[s]$. As consultas fazem \Find$(x)$ e leem o valor na raiz: $O(\alpha(n))$ amortizado.

\resp{uf-tamanho} Indução: uma árvore de altura $h$ tem pelo menos $2^h$ nós. Ao pendurar $T_v$ (a menor) sob a raiz de $T_w$, com $|T_v|\le|T_w|$, a nova altura é $\max\{h_w,h_v+1\}$. Se é $h_w$, vale pela hipótese. Se é $h_v+1$, o novo tamanho é $|T_v|+|T_w|\ge2|T_v|\ge2\cdot2^{h_v}=2^{h_v+1}$. Logo $2^h\le n$.

\resp{uf-ingenua} Faça \Union$(1,2)$, \Union$(2,3)$, \dots, \Union$(n-1,n)$: a raiz do conjunto de $i$ é sempre $i$ (o último elemento a virar raiz é o de $y$), e cada união faz a raiz atual apontar para o novo elemento, formando a cadeia $1\to2\to\cdots\to n$, de altura $n-1$. Sem compressão, $n$ chamadas \Find$(1)$ custam $\Theta(n)$ cada: $\Theta(n^2)$ no total.

\resp{uf-listas} Na primeira rodada, 4 uniões de listas de tamanho 1: 4 atualizações. Na segunda, 2 uniões de listas de tamanho 2, atualizando 2 ponteiros cada: 4. Na terceira, 1 união de listas de tamanho 4: 4. Total: 12. Para $n=2^k$, são $k$ rodadas com $n/2$ atualizações cada: $\frac n2\log_2n$.

\resp{uf-listasprova} O ponteiro de um elemento $x$ só é atualizado quando a lista de $x$ é a menor das duas (ou igual); a nova lista tem pelo menos o dobro do tamanho da lista anterior de $x$. Como nenhuma lista tem mais de $n$ elementos, o ponteiro de $x$ é atualizado no máximo $\log_2n$ vezes. Somando sobre os $n$ elementos, no máximo $n\log_2n$ atualizações.

\resp{uf-semponderada} Una sempre o conjunto grande ao elemento seguinte, anexando a lista grande à unitária: \Union$(2,1)$, \Union$(3,\{1,2\})$, \dots, com a lista do novo elemento como base. A $i$-ésima união atualiza $i$ ponteiros, e o total é $1+2+\cdots+(n-1)=\Theta(n^2)$.

\resp{uf-logstar} Com $F=1,2,4,16,65536,2^{65536}$: $\lgs2=1$, $\lgs16=3$, $\lgs65536=4$ e $\lgs2^{65536}=5$. Como $65536<10^{80}\le2^{65536}$, $\lgs10^{80}=5$.

\resp{uf-comp} Uniões efetivas: $12$, $34$, $23$, $56$, $78$, $89$ (seis). A aresta $65$ liga vértices que já estão juntos (repete $56$) e a aresta $97$ fecha o ciclo $7,8,9$. Componentes finais: $\{1,2,3,4\}$, $\{5,6\}$, $\{7,8,9\}$ e $\{10\}$, ou seja, $10-6=4$ componentes.

\resp{uf-percol} Crie um elemento por casa e dois elementos virtuais, $\top$ e $\bot$. Ao abrir uma casa, una-a a cada vizinha já aberta; se ela está na primeira linha, una-a a $\top$; se está na última, a $\bot$. Há caminho se e somente se \Find$(\top)=$ \Find$(\bot)$. Cada abertura faz no máximo 6 operações; com rank e compressão, $n^2$ aberturas custam $O(n^2\,\alpha(n^2))$.

\resp{uf-equiv} Cada declaração $x\equiv y$ é uma \Union$(x,y)$, e cada consulta compara \Find$(x)$ com \Find$(y)$. No exemplo, $a,b,c,d$ ficam equivalentes ($a\equiv b$, $b\equiv c$, $c\equiv d$), e $e,f$ formam outra classe.

\resp{uf-kruskal} Para cada uma das $m$ arestas, duas operações \Find{} (uma por extremo): $2m$. Uniões efetivas: exatamente $n-1$, uma por aresta da árvore. Com rank e compressão, essas operações custam $O(m\,\alpha(n))$, muito menos que a ordenação das arestas, $O(m\log m)=O(m\log n)$, que domina o algoritmo.

\resp{uf-rankcomp} Atualizar os ranks exigiria recalcular alturas de subárvores depois de cada compressão, o que custaria caro. Com compressão, o rank deixa de ser a altura e passa a ser uma \emph{cota superior} para ela. Ele continua servindo para decidir as uniões porque as propriedades usadas na análise --- uma árvore de rank $k$ tem pelo menos $2^k$ nós e o rank cresce ao subir na árvore --- dependem apenas de como os ranks evoluem nas uniões, e não da altura real.
